\documentclass[11pt,a4paper]{article}

\usepackage[T1]{fontenc}
\usepackage[utf8]{inputenc}
\usepackage{lmodern}
\usepackage[brazil]{babel}
\usepackage[margin=2.5cm]{geometry}
\usepackage{enumitem}
\usepackage{fancyvrb}
\usepackage{titlesec}
\usepackage[hidelinks]{hyperref}

\setlength{\parskip}{0.5em}
\setlength{\parindent}{0pt}
\setlist{itemsep=0.2em, topsep=0.3em}

\newcommand{\codigo}[1]{\texttt{#1}}

\title{Interoperabilidade em Saúde:\\FHIR HL7, RNDS e BR-Core}
\author{RecBitSUS: Recife e UFPE na batida do SUS digital\\ Aula 7}
\date{}

\begin{document}

\maketitle

% ---------------------------------------------------------------
\section{Introdução à interoperabilidade}
% A partir daqui, toda seção começa em página nova
\renewcommand{\sectionbreak}{\clearpage}

\subsection{Definição}

Considere uma viagem organizada por meio de três serviços distintos: a
passagem é comprada no aplicativo da companhia aérea, o hotel é reservado
em outro site e o transporte é solicitado em um terceiro aplicativo. Os
sistemas pertencem a empresas diferentes, mas o horário do voo aparece no
calendário do celular, o hotel recebe o nome do hóspede corretamente e o
motorista localiza o ponto de embarque. Isso é possível porque os sistemas
seguem regras previamente acordadas para a troca de informações.

\textbf{Interoperabilidade} é a capacidade de sistemas diferentes
trocarem dados e utilizarem as informações recebidas sem retrabalho manual.

É importante distinguir dois aspectos dessa definição. Trocar dados
significa que a informação sai de um sistema e chega a outro. Utilizar os
dados significa que o sistema receptor consegue interpretá-los
corretamente: sabe do que se trata cada campo, em que formato está cada
valor e o que cada código representa. Um arquivo enviado por e-mail e
digitado novamente à mão em outro sistema é troca de dados, mas não é
interoperabilidade.

\subsection{Interoperabilidade na saúde}

Considere uma paciente atendida em uma UPA, depois em uma unidade de saúde
da família e, na semana seguinte, em um hospital. Cada serviço utiliza seu
próprio sistema. Na ausência de interoperabilidade:

\begin{itemize}
  \item o histórico precisa ser relatado novamente a cada atendimento;
  \item exames são repetidos porque o resultado anterior não está
    acessível;
  \item alergias e medicamentos em uso podem passar despercebidos;
  \item gestores não conseguem acompanhar o cuidado de forma integrada.
\end{itemize}

Com interoperabilidade, cada sistema envia e recebe os dados da pessoa
atendida, e o profissional dispõe do quadro clínico completo no momento do
atendimento.

Os benefícios se dividem em três níveis:

\begin{itemize}
  \item \textbf{Para o paciente:} menos repetição de exames e de
    relatos, e maior segurança no cuidado, já que informações como
    alergias acompanham a pessoa.
  \item \textbf{Para o profissional:} decisões clínicas tomadas com
    informação completa e atualizada.
  \item \textbf{Para a gestão:} dados consolidados que permitem
    planejar ações e acompanhar a rede de atenção como um todo.
\end{itemize}

\subsection{A necessidade de um padrão}

Para que dois sistemas se comuniquem, a conexão entre eles não é
suficiente. É necessário haver acordo sobre:

\begin{itemize}
  \item quais informações existem (paciente, exame, diagnóstico\ldots);
  \item como cada informação é representada (nome dos campos, formato de
    datas);
  \item o significado de cada código (qual código representa
    ``febre'').
\end{itemize}

Esse acordo constitui o \textbf{padrão de comunicação}.

Sem esse acordo, problemas simples aparecem com facilidade. Uma data
escrita como \codigo{03/04/2026} pode ser lida como 3 de abril ou como 4 de
março, dependendo do sistema. O sexo de uma pessoa pode ser registrado
como ``F'', ``Feminino'' ou ``2''. Um padrão elimina essas ambiguidades ao
definir, de forma única, como cada informação deve ser escrita.

\subsection*{Pontos principais}

\begin{itemize}
  \item Interoperabilidade é trocar dados \emph{e} conseguir utilizá-los.
  \item Na saúde, ela reduz repetição de exames, melhora a segurança do
    paciente e apoia a gestão.
  \item Para sistemas diferentes se entenderem, é preciso um padrão
    comum que defina informações, formatos e códigos.
\end{itemize}

% ---------------------------------------------------------------
\section{Padrões de comunicação e o FHIR}

Duas pessoas que não falam o mesmo idioma podem se comunicar por gestos,
mas o processo é lento e sujeito a erros. Da mesma forma, se cada par de
sistemas definir seu próprio formato de troca, o número de integrações
cresce rapidamente: 5 sistemas exigiriam até 10 integrações distintas e 10
sistemas, até 45.

Com um padrão comum, cada sistema implementa a comunicação uma única vez e
passa a se comunicar com todos os demais. Na saúde digital, um dos padrões
mais utilizados atualmente é o FHIR.

\subsection{O que é o FHIR}

FHIR (\emph{Fast Healthcare Interoperability Resources}) é um padrão de
troca de dados em saúde publicado pela HL7, organização internacional que
desenvolve padrões para informática em saúde.

Toda informação de saúde é dividida em unidades chamadas
\textbf{recursos} (\emph{resources}). Cada recurso possui uma estrutura
definida, e os recursos se relacionam por meio de referências.

Site oficial: \url{https://www.hl7.org/fhir/}

\textbf{Atenção às versões:} o endereço acima apresenta a versão mais
recente da especificação (R5). O Brasil adota a versão R4 (RNDS e
BR-Core). Para consultar a versão utilizada neste material, acesse:
\url{https://hl7.org/fhir/R4/}

\subsection{Características do FHIR}

\begin{itemize}
  \item \textbf{Modular:} a informação é dividida em recursos pequenos e
    independentes, que podem ser combinados conforme a necessidade.
  \item \textbf{Baseado em referências:} um recurso não repete dados de
    outro; ele aponta para ele. Uma observação, por exemplo, referencia o
    paciente em vez de copiar seus dados.
  \item \textbf{Extensível:} quando o padrão não prevê um campo
    necessário, é possível acrescentá-lo por meio de extensões, sem
    quebrar a compatibilidade.
  \item \textbf{Compatível com a web:} o FHIR define uma API baseada em
    HTTP, na qual recursos podem ser lidos, criados, atualizados e
    buscados por meio de endereços (URLs).
\end{itemize}

\subsection{Estrutura de um recurso}

Todo recurso FHIR possui:

\begin{itemize}
  \item um \textbf{tipo} (\codigo{resourceType}), que indica o que ele
    representa: Patient, Observation, Condition, entre outros;
  \item um \textbf{identificador} (\codigo{id}), que o distingue dos
    demais recursos do mesmo tipo;
  \item \textbf{elementos} próprios do tipo, como nome e data de
    nascimento no Patient, ou valor e unidade na Observation.
\end{itemize}

Cada elemento tem um \textbf{tipo de dado}. Alguns tipos que aparecem nos
exemplos deste material:

\begin{description}
  \item[string, date, dateTime] Texto, data (AAAA-MM-DD) e data com hora.
  \item[Coding] Um código de um catálogo: \codigo{system}, \codigo{code}
    e \codigo{display}.
  \item[CodeableConcept] Um conceito representado por um ou mais
    \codigo{Coding}, acompanhados de um texto livre (\codigo{text}).
  \item[Quantity] Um valor numérico com unidade.
  \item[Reference] Uma referência a outro recurso.
\end{description}

\subsection{Recursos principais}

O FHIR define dezenas de recursos. Neste material, são utilizados três:

\begin{description}
  \item[Patient (Paciente)] Dados de identificação: nome, data de
    nascimento, sexo, identificadores (como CPF e CNS).
  \item[Observation (Observação)] Medida ou achado: temperatura,
    pressão arterial, resultado de exame, peso.
  \item[Condition (Condição)] Problema de saúde ou diagnóstico, por
    exemplo, registrado com um código CID.
\end{description}

A estrutura completa de cada recurso (campos, tipos, obrigatoriedade e
diagrama UML) está na página do recurso na especificação oficial da HL7,
que é a representação de referência do modelo:

\begin{itemize}
  \item Patient: \url{https://hl7.org/fhir/R4/patient.html}
  \item Observation: \url{https://hl7.org/fhir/R4/observation.html}
  \item Condition: \url{https://hl7.org/fhir/R4/condition.html}
  \item Bundle: \url{https://hl7.org/fhir/R4/bundle.html}
\end{itemize}

\subsection{Formatos de representação}

Um mesmo recurso pode ser representado em mais de um formato:

\begin{itemize}
  \item \textbf{JSON}: o mais utilizado em APIs modernas;
  \item \textbf{XML}: também suportado pela especificação.
\end{itemize}

O conteúdo é o mesmo; muda apenas a forma de escrevê-lo.

\subsection*{Pontos principais}

\begin{itemize}
  \item O FHIR é um padrão da HL7 que organiza a informação em recursos.
  \item Os recursos são modulares, extensíveis e ligados por referências.
  \item O Brasil utiliza a versão R4.
  \item Patient, Observation e Condition são os recursos usados neste
    material.
\end{itemize}

% ---------------------------------------------------------------
\section{JSON no contexto do FHIR}

\subsection{O que é JSON}

JSON (\emph{JavaScript Object Notation}) é um formato de texto para
organizar dados em pares \codigo{"nome": valor}. Regras básicas:

\begin{itemize}
  \item objetos ficam entre chaves \codigo{\{ \}};
  \item listas ficam entre colchetes \codigo{[ ]};
  \item textos vão entre aspas duplas;
  \item números, \codigo{true}, \codigo{false} e \codigo{null} não levam
    aspas;
  \item itens são separados por vírgula.
\end{itemize}

Exemplo mínimo:

\begin{Verbatim}[frame=single,fontsize=\footnotesize]
{
  "nome": "Maria",
  "idade": 41,
  "alergias": ["dipirona", "penicilina"]
}
\end{Verbatim}

Neste exemplo há um objeto com três pares. O valor de \codigo{"nome"} é um
texto, o de \codigo{"idade"} é um número e o de \codigo{"alergias"} é uma
lista com dois textos.

Erros comuns ao escrever JSON à mão:

\begin{itemize}
  \item esquecer uma vírgula entre dois itens, ou deixar uma vírgula após
    o último item;
  \item usar aspas simples no lugar de aspas duplas;
  \item esquecer de fechar uma chave ou um colchete.
\end{itemize}

Qualquer um desses erros torna o JSON inválido, e o sistema receptor o
rejeita por completo.

\subsection{Por que o FHIR utiliza JSON}

\begin{itemize}
  \item \textbf{Leve:} ocupa pouco espaço e trafega rapidamente.
  \item \textbf{Legível:} pode ser compreendido sem ferramentas
    especiais.
  \item \textbf{Amplo suporte:} todas as linguagens de programação
    leem e escrevem JSON.
\end{itemize}

\subsection{Exemplo: recurso Patient}

Todo recurso FHIR inicia informando o seu tipo, no campo
\codigo{"resourceType"}.

\begin{Verbatim}[frame=single,fontsize=\footnotesize]
{
  "resourceType": "Patient",
  "id": "paciente-001",
  "identifier": [
    {
      "system": "http://www.saude.gov.br/fhir/r4/NamingSystem/cpf",
      "value": "00000000000"
    }
  ],
  "name": [
    {
      "use": "official",
      "family": "Silva",
      "given": ["Maria", "Aparecida"]
    }
  ],
  "gender": "female",
  "birthDate": "1985-03-21"
}
\end{Verbatim}

Interpretação dos campos:

\begin{itemize}
  \item \codigo{"identifier"}: documento que identifica a pessoa. O campo
    \codigo{"system"} indica a qual lista de identificadores o valor
    pertence (neste caso, o CPF). O CPF do exemplo é fictício.
  \item \codigo{"name"}: é uma lista, pois uma pessoa pode ter mais de um
    nome. \codigo{"use"} indica o tipo de uso do nome (por exemplo,
    \codigo{official}), \codigo{"given"} armazena os nomes e
    \codigo{"family"}, o sobrenome.
  \item \codigo{"gender"}: utiliza valores definidos pelo padrão
    (\codigo{female}, \codigo{male}, \codigo{other}, \codigo{unknown}).
  \item \codigo{"birthDate"}: sempre no formato AAAA-MM-DD.
\end{itemize}

Observe que \codigo{"identifier"} e \codigo{"name"} são listas, mesmo
quando há apenas um item. Isso permite que o mesmo recurso armazene, por
exemplo, CPF e CNS, ou um nome oficial e um nome social.

\subsection{Exemplo: recurso Observation}

\begin{Verbatim}[frame=single,fontsize=\footnotesize]
{
  "resourceType": "Observation",
  "id": "obs-001",
  "status": "final",
  "category": [
    {
      "coding": [
        {
          "system": "http://terminology.hl7.org/CodeSystem/observation-category",
          "code": "vital-signs",
          "display": "Vital Signs"
        }
      ]
    }
  ],
  "code": {
    "coding": [
      {
        "system": "http://loinc.org",
        "code": "8310-5",
        "display": "Body temperature"
      }
    ],
    "text": "Temperatura corporal"
  },
  "subject": {
    "reference": "Patient/paciente-001"
  },
  "effectiveDateTime": "2026-10-09T09:30:00-03:00",
  "valueQuantity": {
    "value": 38.5,
    "unit": "°C",
    "system": "http://unitsofmeasure.org",
    "code": "Cel"
  }
}
\end{Verbatim}

Interpretação dos campos:

\begin{itemize}
  \item \codigo{"status"}: situação do registro. Alguns valores previstos
    pelo padrão são \codigo{registered} (registrado), \codigo{preliminary}
    (preliminar), \codigo{final} (finalizado) e \codigo{amended}
    (corrigido).
  \item \codigo{"category"}: agrupa a observação em uma categoria, como
    sinais vitais (\codigo{vital-signs}).
  \item \codigo{"code"}: o que foi medido. O código vem do LOINC, catálogo
    internacional de exames e medidas. O campo \codigo{"text"} traz a
    descrição legível.
  \item \codigo{"subject"}: a quem pertence a observação. A referência
    aponta para o Patient.
  \item \codigo{"effectiveDateTime"}: quando a medida foi realizada. O
    trecho \codigo{-03:00} indica o fuso horário (horário de Brasília).
  \item \codigo{"valueQuantity"}: valor, unidade e código padronizado da
    unidade. O código \codigo{Cel} pertence ao UCUM, padrão internacional
    de unidades de medida, e representa graus Celsius.
\end{itemize}

\subsection{Exemplo: recurso Condition com CID}

\begin{Verbatim}[frame=single,fontsize=\footnotesize]
{
  "resourceType": "Condition",
  "id": "cond-001",
  "clinicalStatus": {
    "coding": [
      {
        "system": "http://terminology.hl7.org/CodeSystem/condition-clinical",
        "code": "active"
      }
    ]
  },
  "code": {
    "coding": [
      {
        "system": "http://www.saude.gov.br/fhir/r4/CodeSystem/BRCID10",
        "code": "R50.9",
        "display": "Febre não especificada"
      }
    ],
    "text": "Febre"
  },
  "subject": {
    "reference": "Patient/paciente-001"
  },
  "recordedDate": "2026-10-09"
}
\end{Verbatim}

\textbf{O que é o CID:} a Classificação Internacional de Doenças, da
Organização Mundial da Saúde, é o catálogo que atribui um código a cada
doença ou problema de saúde. Utiliza-se aqui a versão CID-10 (R50.9 =
febre não especificada).

\textbf{CID-10 e CIAP-2:} o CID-10 é o catálogo usado para registrar
\textbf{diagnósticos}. Para registrar \textbf{problemas} ou motivos de
atendimento, sobretudo na atenção primária, utiliza-se também a CIAP-2
(Classificação Internacional de Atenção Primária). Ambos podem alimentar o
campo \codigo{"code"} da Condition; o que muda é o \codigo{"system"}, que
indica o catálogo de origem.

Interpretação dos campos:

\begin{itemize}
  \item \codigo{"clinicalStatus"}: situação clínica do problema.
    \codigo{active} indica que o problema está ativo. O padrão prevê
    outros valores, como \codigo{inactive}, \codigo{remission} e
    \codigo{resolved}.
  \item \codigo{"code"}: qual é o problema de saúde, expresso por um
    código do CID-10.
  \item \codigo{"recordedDate"}: data em que o problema foi registrado.
\end{itemize}

\subsection{Como o FHIR representa códigos}

Nos três exemplos, os campos \codigo{"code"} e \codigo{"category"}
seguem a mesma estrutura, chamada \codigo{CodeableConcept}. Ela contém:

\begin{itemize}
  \item uma lista \codigo{"coding"}, na qual cada item é um código de um
    catálogo;
  \item um campo \codigo{"text"}, com a descrição em texto livre.
\end{itemize}

Cada item de \codigo{"coding"} tem três elementos: \codigo{"system"}
(catálogo de origem do código), \codigo{"code"} (o código em si) e
\codigo{"display"} (o nome legível). Sem o \codigo{"system"}, um código
isolado é ambíguo, pois o mesmo número pode ter significados diferentes em
catálogos diferentes.

A lista \codigo{"coding"} permite que o mesmo conceito seja expresso em
mais de um catálogo ao mesmo tempo, e o campo \codigo{"text"} preserva o
que o profissional escreveu, caso não exista código adequado.

\subsection*{Pontos principais}

\begin{itemize}
  \item JSON organiza dados em objetos (\codigo{\{ \}}) e listas
    (\codigo{[ ]}).
  \item Todo recurso FHIR começa com \codigo{"resourceType"}.
  \item Códigos são sempre acompanhados de seu \codigo{"system"}.
  \item Recursos apontam para o paciente por meio de
    \codigo{"subject"}.
\end{itemize}

% ---------------------------------------------------------------
\section{Bundles}

\subsection{Conceito}

Um Bundle é um pacote que agrupa vários recursos em um único JSON. É
utilizado, por exemplo, para enviar de uma só vez o paciente e tudo o que
foi registrado sobre ele em um atendimento.

Principais tipos de Bundle: \codigo{collection} (coleção de recursos),
\codigo{document} (documento clínico), \codigo{transaction} (conjunto de
operações que devem ser aplicadas em conjunto) e \codigo{searchset}
(resultado de uma busca). Os exemplos deste material utilizam o tipo
\codigo{collection}.

Enviar os recursos agrupados traz duas vantagens: reduz o número de
requisições entre os sistemas e mantém juntas as informações que pertencem
ao mesmo contexto, como as de um mesmo atendimento.

\subsection{Exemplo: paciente, observação e condição}

{\scriptsize
\begin{Verbatim}[frame=single]
{
  "resourceType": "Bundle",
  "type": "collection",
  "entry": [
    {
      "fullUrl": "urn:uuid:11111111-1111-1111-1111-111111111111",
      "resource": {
        "resourceType": "Patient",
        "identifier": [
          {
            "system": "http://www.saude.gov.br/fhir/r4/NamingSystem/cpf",
            "value": "00000000000"
          }
        ],
        "name": [
          { "family": "Silva", "given": ["Maria", "Aparecida"] }
        ],
        "gender": "female",
        "birthDate": "1985-03-21"
      }
    },
    {
      "fullUrl": "urn:uuid:22222222-2222-2222-2222-222222222222",
      "resource": {
        "resourceType": "Observation",
        "status": "final",
        "code": {
          "coding": [
            {
              "system": "http://loinc.org",
              "code": "8310-5",
              "display": "Body temperature"
            }
          ],
          "text": "Temperatura corporal"
        },
        "subject": {
          "reference": "urn:uuid:11111111-1111-1111-1111-111111111111"
        },
        "effectiveDateTime": "2026-10-09T09:30:00-03:00",
        "valueQuantity": {
          "value": 38.5,
          "unit": "°C",
          "system": "http://unitsofmeasure.org",
          "code": "Cel"
        }
      }
    },
    {
      "fullUrl": "urn:uuid:33333333-3333-3333-3333-333333333333",
      "resource": {
        "resourceType": "Condition",
        "clinicalStatus": {
          "coding": [
            {
              "system": "http://terminology.hl7.org/CodeSystem/condition-clinical",
              "code": "active"
            }
          ]
        },
        "code": {
          "coding": [
            {
              "system": "http://www.saude.gov.br/fhir/r4/CodeSystem/BRCID10",
              "code": "R50.9",
              "display": "Febre não especificada"
            }
          ],
          "text": "Febre"
        },
        "subject": {
          "reference": "urn:uuid:11111111-1111-1111-1111-111111111111"
        },
        "recordedDate": "2026-10-09"
      }
    }
  ]
}
\end{Verbatim}
}

Cada item de \codigo{"entry"} contém um recurso e um \codigo{"fullUrl"},
que funciona como o endereço do recurso dentro do pacote. A Observation e a
Condition utilizam esse endereço em \codigo{"subject.reference"} para
apontar para o mesmo Patient. É dessa forma que os recursos se relacionam
dentro do Bundle.

\subsection{Referências dentro e fora do Bundle}

Nos exemplos individuais da seção anterior, a referência ao paciente foi
escrita como \codigo{"Patient/paciente-001"}, ou seja, tipo do recurso e
\codigo{id}. Esse formato pressupõe que o Patient já exista em algum
servidor com esse identificador.

Dentro de um Bundle, o paciente ainda pode não ter um identificador
definitivo. Por isso, cada entrada recebe um \codigo{"fullUrl"} temporário
no formato \codigo{urn:uuid:...}, e as referências apontam para ele. Um UUID
é um identificador gerado de forma que seja praticamente único, o que evita
conflito entre recursos.

A regra prática é: o valor de \codigo{"reference"} deve coincidir com o
\codigo{"fullUrl"} de alguma entrada do pacote. Se não coincidir, a
referência fica sem destino.

\subsection*{Pontos principais}

\begin{itemize}
  \item O Bundle agrupa recursos relacionados em um único envio.
  \item Cada entrada tem um \codigo{"fullUrl"}, que serve de endereço
    dentro do pacote.
  \item As referências entre recursos devem apontar para um
    \codigo{"fullUrl"} existente.
\end{itemize}

% ---------------------------------------------------------------
\section{RNDS e BR-Core}

\subsection{RNDS -- Rede Nacional de Dados em Saúde}

A RNDS é a plataforma nacional do Ministério da Saúde para integração de
dados em saúde. Seu objetivo é permitir que estabelecimentos públicos e
privados enviem e consultem informações de forma padronizada, compondo o
histórico clínico do cidadão, que pode ser acessado, por exemplo, por meio
do Conecte SUS.

A RNDS padroniza em âmbito nacional a linguagem de troca de informações em
saúde, e essa linguagem é o FHIR (versão R4). Os dados enviados à RNDS são
validados com base em modelos publicados pelo Ministério da Saúde.

Esses modelos de dados, com os campos de cada tipo de registro, estão no
guia de integração da RNDS: \url{https://rnds-guia.saude.gov.br/}

Na prática, isso significa que um sistema de um município, de um hospital
ou de uma clínica privada, desenvolvido por qualquer fornecedor, pode
trocar informações com a rede nacional desde que adote o mesmo padrão.

\subsection{BR-Core}

O BR-Core é o conjunto de definições que adapta o FHIR à realidade
brasileira. O FHIR é um padrão internacional e deixa diversas decisões em
aberto, como quais identificadores utilizar e quais catálogos de códigos
adotar. O BR-Core define essas escolhas para o Brasil, por exemplo:

\begin{itemize}
  \item identificadores nacionais, como CPF e CNS;
  \item catálogos de códigos brasileiros (como o CID-10 usado no exemplo);
  \item perfis: regras que definem quais campos de um recurso são
    obrigatórios.
\end{itemize}

Em termos de analogia, se o FHIR é o idioma, o BR-Core corresponde às
regras de uso e ao vocabulário acordados para o país.

\subsection{Perfis e validação}

Um \textbf{perfil} é um conjunto de restrições aplicado a um recurso. O
recurso Patient do FHIR, por exemplo, deixa quase todos os campos
opcionais. Um perfil pode estabelecer que determinados campos passam a ser
obrigatórios, que certos identificadores devem ser usados e que
determinados catálogos de códigos devem ser respeitados.

A \textbf{validação} é a verificação de que um recurso respeita as regras
do perfil. Um recurso pode ser um JSON válido e, ainda assim, falhar na
validação por não cumprir o perfil exigido, por exemplo por não conter um
campo obrigatório.

Por isso, existem três níveis de correção que convém distinguir:

\begin{enumerate}
  \item o JSON está bem formado (sintaxe correta);
  \item o recurso respeita a estrutura do FHIR (campos e tipos corretos);
  \item o recurso respeita o perfil exigido (por exemplo, o do BR-Core).
\end{enumerate}

\textbf{Observação:} nos exemplos deste material, o \codigo{"system"} do
CPF e do CID segue o padrão brasileiro. Em um envio real à RNDS, é
necessário conferir as URLs e os perfis na documentação oficial vigente.

\subsection*{Pontos principais}

\begin{itemize}
  \item A RNDS é a rede nacional que utiliza o FHIR R4 como linguagem de
    troca.
  \item O BR-Core adapta o FHIR ao Brasil: identificadores, catálogos e
    perfis.
  \item Perfis definem regras adicionais, e a validação verifica se elas
    foram cumpridas.
\end{itemize}

% ---------------------------------------------------------------
\section{Atividade prática}

\textbf{Objetivo:} observar a geração de um Bundle FHIR e alterar o JSON
para compreender como a estrutura reage às modificações.

\textbf{Repositório:}
\url{https://github.com/Daniel-Nas/Interoperabilidade-RecBitSus}

\textbf{Procedimento:}

\begin{enumerate}
  \item Clonar o repositório.
  \item Seguir o README para instalar as dependências e executar o
    projeto.
  \item Preencher o formulário e gerar um Bundle.
  \item Abrir o JSON gerado e identificar: o Patient, a Observation e a
    Condition; o \codigo{"fullUrl"} de cada recurso; as referências em
    \codigo{"subject"}.
\end{enumerate}

Caso prefira, acesse \textbf{[LINK DO DEPLOY NO VERCEL]}

\textbf{Teste em um servidor FHIR (opcional):} o JSON gerado também pode
ser enviado a um servidor FHIR de testes. O HAPI FHIR, implementação
de referência muito utilizada, mantém um servidor público R4 para
experimentação: \url{https://hapi.fhir.org/baseR4}. Nele, é possível
enviar cada recurso (por exemplo, \codigo{POST} em
\codigo{https://hapi.fhir.org/baseR4/Patient}) e depois consultá-lo por
\codigo{GET}. O envio pode ser feito sem programar nada, apenas com o JSON
dos exemplos desta apostila colado como corpo da requisição em uma
ferramenta como o Postman. Por ser um servidor público, \textbf{não envie
dados reais de pessoas}.

\textbf{Modificações propostas:}

\begin{enumerate}[label=\alph*)]
  \item Trocar o CID do exemplo (por exemplo, R50.9 por outro código) e
    alterar o \codigo{"display"} correspondente.
  \item Alterar o valor da temperatura e a data da observação.
  \item Adicionar um segundo nome em \codigo{"given"} no Patient.
  \item Apagar o \codigo{"system"} de um \codigo{"coding"} e discutir o
    que se perde.
  \item Remover um campo obrigatório e observar o resultado.
\end{enumerate}

\textbf{Questões para discussão:}

\begin{itemize}
  \item Por que o \codigo{"system"} é tão importante quanto o
    \codigo{"code"}?
  \item O que acontece com um Bundle quando a referência de
    \codigo{"subject"} aponta para um recurso inexistente?
  \item Que informações do cotidiano de uma unidade de saúde ainda não
    podem ser representadas apenas com Patient, Observation e Condition?
\end{itemize}

% ---------------------------------------------------------------
\section*{Glossário}

\begin{description}
  \item[BR-Core] Adaptação do FHIR para o Brasil.
  \item[Bundle] Pacote que agrupa vários recursos FHIR.
  \item[CID] Classificação Internacional de Doenças (OMS).
  \item[CIAP-2] Classificação Internacional de Atenção Primária.
  \item[CNS] Cartão Nacional de Saúde.
  \item[CodeableConcept] Tipo de dado do FHIR que reúne um ou mais códigos
    e um texto livre.
  \item[Condition] Recurso FHIR para problema de saúde ou diagnóstico.
  \item[FHIR] Padrão de troca de dados em saúde criado pela HL7.
  \item[HAPI FHIR] Implementação de referência do FHIR em Java, com
    servidor público de testes.
  \item[HL7] Organização internacional de padrões em saúde.
  \item[Interoperabilidade] Capacidade de sistemas trocarem e usarem dados.
  \item[JSON] Formato de texto para estruturar dados.
  \item[LOINC] Catálogo internacional de exames e medidas.
  \item[Observation] Recurso FHIR para medidas e achados.
  \item[Patient] Recurso FHIR para dados do paciente.
  \item[Perfil] Conjunto de restrições aplicado a um recurso FHIR.
  \item[Recurso] Unidade básica de informação do FHIR.
  \item[RNDS] Rede Nacional de Dados em Saúde.
  \item[UCUM] Padrão de unidades de medida.
  \item[Validação] Verificação de que um recurso respeita as regras de um
    perfil.
\end{description}

% ---------------------------------------------------------------
% Desativa a quebra de página apenas para Referências
\renewcommand{\sectionbreak}{}
\section*{Referências}

\begin{itemize}
  \item HL7 FHIR (R4): \url{https://hl7.org/fhir/R4/}
  \item HL7 FHIR (versão atual): \url{https://www.hl7.org/fhir/}
  \item RNDS, guia de integração: \url{https://rnds-guia.saude.gov.br/}
  \item HAPI FHIR, servidor de testes R4: \url{https://hapi.fhir.org/baseR4}
  \item Repositório da atividade prática:
    \url{https://github.com/Daniel-Nas/Interoperabilidade-RecBitSus}
  \item iPES: \url{https://ipes.tech/}
  \item LOINC: \url{https://loinc.org/}
\end{itemize}

\end{document}